% ===========================================================================
\section{Discussion}
\label{sec:discussion}
% ===========================================================================

\subsection{What anchoring buys, and what it costs}

Anchoring gives the sector total three properties that a bottom-up estimate does not have. It
agrees with the national fuel statistics exactly rather than approximately. It is insulated
from the fleet model, so the least supported input in the sector cannot move the most used
output. And it makes the boundary problem---which customer classes are road, and what a
retail sales channel can and cannot see---visible as an executable ledger rather than
absorbed into an unexamined assumption about what fraction of diesel is "road".

The cost is real and worth naming. The class composition in Section~\ref{sec:results} is
determined entirely by the fleet model, with no independent check, so a reader wanting
emissions by vehicle type gets a number whose uncertainty this paper cannot quantify. A
bottom-up inventory would produce the same weak composition but would additionally propagate
it into the total; anchoring confines it. That is a better trade, not a free one.

There is also a subtler point, which Supplementary Note~S1 quantifies. Removing the fleet
uncertainty from the total leaves the shared fuel-carbon factor as the dominant term. An
earlier edition read the cross-sector correlation scenario as roughly doubling the band;
the common-random-numbers rebuild showed that spread to be a sampling artifact, and both
treatments of the pooled Monte Carlo now print the same $\pm\SI{12.9}{\percent}$, beside the
sector's own Approach~1 interval of $\pm\SI{11.4}{\percent}$ that the results report as the
number to quote. Anchoring does not make the
sector more certain; it makes the remaining uncertainty legible.

\subsection{What the re-anchor changed, and what it teaches}

Reading the controls from their primary table moved the sector by six megatonnes and
flipped the EDGAR comparison from above the benchmark to below it. Neither number was
the point. The previous edition's control was arithmetically inconsistent with its own
source note, and the matched-year benchmark agreement it produced was coincidence---the
wrong annualisation happened to land near the right vintage. The lesson generalises: a
benchmark comparison tests consistency, not provenance, and an inventory that treats
agreement with EDGAR as validation has its evidentiary arrow backwards.

The re-anchor also shrank the boundary problem rather than merely renaming it. Bulk
industrial, utility, government and railway diesel left the road control by the structure
of the source; what remains unresolved is the retail channel's end-use split, bounded by
the ten-year road-share panel and carried as status rows. The residual ask of Pakistani
statistics is now precise: not "publish a sectoral split"---OCAC does---but an end-use
survey of the retail diesel channel, separating road vehicles from tube wells and captive
gensets at the pump.

The asymmetry the previous edition claimed survives, strengthened: this inventory can name
what is wrong with its own transport figure more precisely than any external comparison
can. A ratio says two numbers differ; the ledger says which claimant is unresolved, in
which channel, inside which row, within what bound.

\subsection{Implications}

\textbf{Fuel controls and fleet assumptions constrain different outputs.} Fuel anchoring constrains combustion CO$_2$ and isolates the effect of class partition on that total. PM, NO$_x$, evaporation and wear remain sensitive to vehicle technologies, controls, maintenance and operating conditions. A GHG variance screen therefore cannot establish the value of fleet evidence for the air-pollutant product.

\textbf{Retail end-use evidence would constrain the source boundary.} The unresolved fraction of retail diesel used by agricultural and captive equipment affects the road/off-road assignment directly. This evidence serves a different purpose from fleet technology and vehicle-kilometre information, which constrain pollutant factors and class composition. Their value cannot be ranked using fuel-derived CO$_2$ alone.

\textbf{Two-wheelers are the largest single class, but not the largest claim.} Under the
reference allocation they receive \SI{34.1}{\percent} of on-road \COtwo{}; the
superseded run's composition envelope spanned twelve points around it, and in
every one of them heavy goods vehicles and buses together
(\SI{14.02}{\mega\tonne} at the reference) exceed them. The previous edition, splitting
fuel by distance alone, put two-wheelers at \SI{42.3}{\percent} and above heavy diesel;
that ordering was an artefact of the weighting and this edition retires it.

We draw no policy conclusion from the ordering, and the earlier draft's remark that
car-focused electrification addresses only a minority of the sector is withdrawn---not
because it is necessarily wrong, but because nothing here supports it. A ranking whose
inputs are registration counts, assumed kilometres and declared intensities, none of them
measured in Pakistan, is not a basis for allocating policy attention between vehicle
classes. What the composition can support is a weaker and more useful statement: whichever
end of the envelope holds, this is a fleet in which two-wheelers and heavy diesel each
account for a substantial share, and an inventory that resolved only one of them would
mislead.
